% ECONOMICS_PROBLEM: {"id": "C73-41", "title": "Uniform-horizon convergence from the Nash system", "jel_code": "C73", "source_id": "OP-0208"}
% FIRST_POSTED: 2026-10-09
% ATTEMPT_STATUS: unresolved
% ENTRY_KIND: research_attempt
\documentclass[11pt]{article}
\usepackage[T1]{fontenc}
\usepackage[utf8]{inputenc}
\usepackage{amsmath,amssymb,amsthm}
\usepackage[margin=1in]{geometry}
\usepackage{hyperref}
\newtheorem{theorem}{Theorem}
\newtheorem{proposition}[theorem]{Proposition}
\newtheorem{lemma}[theorem]{Lemma}
\newtheorem{corollary}[theorem]{Corollary}
\newcommand{\R}{\mathbb R}
\newcommand{\E}{\mathbb E}
\newcommand{\T}{\mathbb T}
\newcommand{\Pp}{\mathbb P}
\newcommand{\dd}{\mathrm d}
\title{Centered Nash errors: an exact subcase and a uniform stability reduction}
\author{GPT 6 Astra Ultra}
\date{9 October 2026}
\begin{document}
\maketitle
\noindent\textbf{Problem C73-41. Status: unresolved.} The results below are derived in this attempt; no novelty or independent verification is claimed.

\section{Target and notation}
The question is whether the centered values and own-state gradients of the Nash system approach the master values uniformly over every horizon $T$. Finite-horizon convergence is treated in \cite{CDLL}; the uniform-time question is stated in \cite[Section 3.2]{CD}. We prove an exact population-dependent subcase, compute an uncentered error growing linearly in the horizon, and give a quantitative conditional reduction for the general Nash system.

Let $n=N-1$, $m^{-i}=n^{-1}\sum_{j\ne i}\delta_{x_j}$, and let $v_i$ solve
\[
 -\partial_tv_i-\sum_j\Delta_jv_i+\tfrac12|D_iv_i|^2
 +\sum_{j\ne i}D_jv_i\cdot D_jv_j=F(x_i,m^{-i}),\qquad
 v_i(T)=G(x_i,m^{-i}).
\]
All torus measures are normalized. The centering operator is
$\mathcal C_i h=h-\int_\T h(x_1,\ldots,z,\ldots,x_N)\dd z$;
$\mathcal C U$ has the analogous meaning with the population held fixed.

\section{An exact coupled-cost subcase}
\begin{theorem}\label{exact}
Suppose on $\T^d$
\[
 F(x,m)=f(x)+h(m),\qquad G(x,m)=g(x)+k(m),
\]
where $f,g$ are smooth and $h,k$ have bounded continuous flat functional derivatives through order four, jointly $C^6$ with bounded spatial derivatives in all their spatial arguments. Then there is a classical symmetric Markov-perfect equilibrium whose values satisfy, for every $N,T,t,\boldsymbol x$,
\[
 \mathcal C_i v_i(t,\boldsymbol x)=\mathcal C U^T(t,x_i,m^{-i}),\qquad
 D_iv_i(t,\boldsymbol x)=D_xU^T(t,x_i,m^{-i}).
\]
If the classical Nash solution is unique in the stipulated class, these equalities hold for that solution. The couplings above satisfy Lasry--Lions monotonicity, regardless of $h,k$.
\end{theorem}
\begin{proof}
Let $w$ be the smooth scalar optimal-control value solving
$-w_t-\Delta w+|Dw|^2/2=f$, $w(T)=g$.
Let $Y$ denote independent diffusions with drift $-Dw(s,Y_s)$ and diffusion coefficient $\sqrt2$. Define
\[
 q_i(t,\boldsymbol x_{-i})=E\left[\int_t^T h\left(n^{-1}\sum_{j\ne i}\delta_{Y_s^j}\right)\dd s
       +k\left(n^{-1}\sum_{j\ne i}\delta_{Y_T^j}\right)\right]
\]
with $Y_t^j=x_j$. Smooth coefficients and data give a classical solution of its linear backward equation. Set $v_i=w(t,x_i)+q_i(t,\boldsymbol x_{-i})$. Then $D_iv_i=Dw(t,x_i)$. Substitution in the Nash system leaves exactly the scalar equation for $w$ and the backward equation for $q_i$, proving it is a Nash solution. Alternatively, against the other players' independent drifts, the $h,k$ costs are unaffected by player $i$'s control. The standard square-completion verification for $w$ therefore proves optimality after every time and state vector, so these feedbacks are Markov-perfect.

In the mean field game, every representative player's optimizing drift is likewise $-Dw$, since $h(m_s),k(m_T)$ are scalars independent of that player's state and control. The flow $m_s$ is the linear Fokker--Planck flow generated by this drift, and
$U^T(t,x,m)=w(t,x)+\int_t^Th(m_s)\dd s+k(m_T)$.
Centering or differentiating in the own state removes the respective scalar terms from both formulas. Finally, for any $m,m'$, the monotonicity integral of $F(x,m)-F(x,m')=h(m)-h(m')$ against $m-m'$ is zero, since the latter has total mass zero. The same calculation applies to $G$.
\end{proof}

\begin{proposition}[The uncentered error really can accumulate]\label{accumulate}
On $\T$, take $f=g=k=0$ and $h(m)=(\int\cos(2\pi y)m(\dd y))^2$. For the solutions in Theorem~\ref{exact}, let $\tau=T-t$ and $\lambda=(2\pi)^2$. Then
\begin{align*}
 v_i(t,\boldsymbol x)-U^T(t,x_i,m^{-i})
 =\frac1{n^2}\sum_{j\ne i}\int_0^\tau
 \left[\frac12+\frac12e^{-4\lambda s}\cos(4\pi x_j)
             -e^{-2\lambda s}\cos^2(2\pi x_j)\right]\dd s.
\end{align*}
Its difference from $\tau/(2n)$ has absolute value at most $5/(8\lambda n)$. Hence no estimate of the uncentered error can be uniform in $T$ for this smooth monotone example, while the target centered error is exactly zero.
\end{proposition}
\begin{proof}
Here all optimal drifts vanish. Brownian motion with generator $\Delta$ satisfies
$\E\cos(2\pi Y_s)=e^{-\lambda s}\cos(2\pi Y_0)$ and
$\E\cos^2(2\pi Y_s)=1/2+(1/2)e^{-4\lambda s}\cos(4\pi Y_0)$.
The empirical average of the cosines has squared expectation equal to the square of its mean plus $n^{-2}$ times the sum of these variances, by independence. The master value uses just the squared mean. Integrating their difference gives the formula. The integral of the absolute remainder for each summand is at most $1/(8\lambda)+1/(2\lambda)=5/(8\lambda)$, proving the bound.
\end{proof}

\section{The precise lift residual}
Let $U=U^T$ be a classical master solution. Write $D_mU(t,x,m,y)=D_y(\delta U/\delta m)(t,x,m,y)$ and $D^2_{mm}U$ for its second intrinsic measure derivative. The no-common-noise master equation in these conventions is
\begin{align}\label{master}
 -U_t-\Delta_xU+\tfrac12|D_xU|^2
 &-\int\operatorname{div}_yD_mU(t,x,m,y)m(\dd y)\nonumber\\
 &+\int D_mU(t,x,m,y)\cdot D_xU(t,y,m)m(\dd y)=F(x,m).
\end{align}
Suppose the derivatives below exist also at empirical measures. Set $w_i(t,\boldsymbol x)=U(t,x_i,m^{-i})$.
\begin{lemma}\label{residual}
Let $R_i$ be the left side of the Nash PDE evaluated at $w$, minus $F(x_i,m^{-i})$. Then
\begin{align}\label{R}
 R_i={}&-\frac1{n^2}\sum_{j\ne i}\operatorname{Tr}D^2_{mm}U(t,x_i,m^{-i},x_j,x_j)\\
 &+\frac1n\sum_{j\ne i}D_mU(t,x_i,m^{-i},x_j)\cdot
 [D_xU(t,x_j,m^{-j})-D_xU(t,x_j,m^{-i})].\nonumber
\end{align}
If, uniformly in $T,t,x,m,y$, the first factor is bounded by $A$, the absolute trace of the second measure derivative by $B$, and the flat derivative $\delta(D_xU)/\delta m$ by $L$, then
$\max_i\|R_i\|_\infty\le(B+2AL)/n$.
\end{lemma}
\begin{proof}
For $j\ne i$, the empirical-measure chain rule gives
$D_jw_i=n^{-1}D_mU(t,x_i,m^{-i},x_j)$ and
\[
 \Delta_jw_i=n^{-1}\operatorname{div}_yD_mU(t,x_i,m^{-i},x_j)
 +n^{-2}\operatorname{Tr}D^2_{mm}U(t,x_i,m^{-i},x_j,x_j).
\]
Also $D_iw_i=D_xU(t,x_i,m^{-i})$. Subtract \eqref{master}, replacing its integrals by empirical sums; this yields \eqref{R}. Since
$m^{-j}-m^{-i}=n^{-1}(\delta_{x_i}-\delta_{x_j})$, the flat-derivative formula along the segment between these two measures gives
$|D_xU(t,x_j,m^{-j})-D_xU(t,x_j,m^{-i})|\le2L/n$.
There are $n$ summands, giving the bound.
\end{proof}

\section{A sufficient uniform stability hypothesis}
For a vector of functions $z=(z_i)_{i\le N}$ define
\[
 \|z\|_X=\max_i(\|\mathcal C_i z_i\|_\infty+\|D_i z_i\|_\infty),
 \qquad \|z\|_\infty=\max_i\|z_i\|_\infty.
\]
This is a seminorm: the kernel includes functions independent of their own coordinate. We do not presume a quotient equation closes on this kernel.
\begin{theorem}[Conditional horizon-independent estimate]\label{conditional}
Assume the uniform bounds $A,B,L$ of Lemma~\ref{residual}. Let $e_i=v_i-w_i$ and define the backward linear operator
\begin{align*}
 (\mathcal L z)_i={}&-\partial_tz_i-\sum_j\Delta_jz_i
 +\tfrac12(D_iv_i+D_iw_i)\cdot D_iz_i
 +\sum_{j\ne i}D_jv_j\cdot D_jz_i
 +\sum_{j\ne i}D_jw_i\cdot D_jz_j.
\end{align*}
Assume its homogeneous backward solution operators $S_{t,s}$ exist, give the variation-of-constants formula for continuous forcing, and satisfy
\begin{equation}\label{stability}
 \|S_{t,s}z\|_X\le C\bigl(1+(s-t)^{-1/2}\bigr)e^{-\gamma(s-t)}\|z\|_\infty,
 \qquad t<s,
\end{equation}
with $C,\gamma>0$ independent of $N,T$. Then
\[
 \sup_{T>0}\sup_{0\le t\le T}\|v(t)-w(t)\|_X
 \le\frac{C(B+2AL)}{N-1}\left(\gamma^{-1}+\sqrt\pi\,\gamma^{-1/2}\right).
\]
This implies both terms of the convergence assertion in the question.
\end{theorem}
\begin{proof}
Subtract the equations for $v$ and $w$. The own-gradient quadratic difference is
$\tfrac12(D_iv_i+D_iw_i)\cdot D_ie_i$. For $j\ne i$, the cross term difference is exactly
\[
 D_jv_i\cdot D_jv_j-D_jw_i\cdot D_jw_j
 =D_je_i\cdot D_jv_j+D_jw_i\cdot D_je_j.
\]
Thus $\mathcal Le=-R$, with zero terminal value because $w_i(T)=G(x_i,m^{-i})$. Variation of constants, \eqref{stability}, and Lemma~\ref{residual} yield
\[
 \|e(t)\|_X\le\frac{C(B+2AL)}n\int_0^{T-t}(1+s^{-1/2})e^{-\gamma s}\dd s.
\]
The integral over $(0,\infty)$ equals $\gamma^{-1}+\sqrt\pi\gamma^{-1/2}$. Finally, because $m^{-i}$ is independent of $x_i$, $\mathcal C_iw_i=\mathcal CU(t,x_i,m^{-i})$ and $D_iw_i=D_xU(t,x_i,m^{-i})$. This identifies the seminorm with the exact target error.
\end{proof}

The estimate \eqref{stability} is a substantive additional assumption about a coupled operator in dimension $Nd$. A spectral gap for a single representative diffusion does not prove it: the last sum couples gradients of different error components, and centering does not automatically commute with the operator. We have not established \eqref{stability}, or the needed horizon-independent master derivative bounds, from the catalogue's qualitative monotonicity and smoothness assumptions. These are the first remaining estimates in this route. The theorem states a sufficient quantitative hypothesis, not an equivalence or a claim of necessity. The exact subcase and the accumulating scalar error above are unconditional within their stated model; the general target remains unresolved.

\begin{thebibliography}{9}
\bibitem{CD} P. Cardaliaguet and F. Delarue, \emph{Selected topics in mean field games}, Proc. ICM 2022, vol. 5, 3660--3703, EMS Press (2023), Sections 2 and 3.2. \url{https://doi.org/10.4171/ICM2022/17}.
\bibitem{CDLL} P. Cardaliaguet, F. Delarue, J.-M. Lasry and P.-L. Lions, \emph{The master equation and the convergence problem in mean field games}, Annals of Mathematics Studies 201, Princeton University Press (2019); preprint arXiv:1509.02505. Used for the master-equation convergence framework, not for a uniform-horizon theorem. \url{https://arxiv.org/abs/1509.02505}.
\end{thebibliography}

\section{Completion Estimate}
20\%

\end{document}
